% Źródła: Eves s. 311--314 (§7.8); Greenberg 2010, s. 208; Wikipedia: Johannes Hjelmslev.
% https://en.wikipedia.org/wiki/Johannes_Hjelmslev

\section{Hjelmslev i odwzorowanie na koło} % lang-pl
\section{Hjelmslev e la rappresentazione nel cerchio} % lang-it
\label{sekcja_hjelmslev}

Model w kole da się otrzymać także na inny sposób: bez rzutowań, samym odwzorowaniem płaszczyzny Łobaczewskiego. % lang-pl
W tej sekcji poznamy duńskiego geometrę, którego twierdzenie o środkach odcinków to odwzorowanie umożliwi, i zobaczymy, jak ta sama konstrukcja działa w geometrii euklidesowej. % lang-pl
Un modello nel cerchio si può ottenere anche in un altro modo: senza proiezioni, con una semplice trasformazione del piano di Lobachevskij. % lang-it
In questa sezione conosceremo il geometra danese il cui teorema sui punti medi renderà possibile questa trasformazione, e vedremo come la stessa costruzione funzioni nella geometria euclidea. % lang-it

Johannes Trolle Hjelmslev urodzi się 7 kwietnia 1873 roku w Hørning w Danii i umrze 16 lutego 1950 roku; będzie profesorem, a później rektorem Uniwersytetu Kopenhaskiego. % lang-pl
Urodzi się jako Johannes Trolle Petersen, a nazwisko Hjelmslev przyjmie, żeby odróżnić się od Juliusa Petersena -- pod dawnym nazwiskiem zostanie w matematyce m.in. twierdzenie Petersena--Morleya. % lang-pl
Będzie ojcem językoznawcy Louisa Hjelmsleva. % lang-pl
Johannes Trolle Hjelmslev nascerà il 7 aprile 1873 a Hørning, in Danimarca, e morirà il 16 febbraio 1950; sarà professore e poi rettore dell'Università di Copenaghen. % lang-it
Nascerà come Johannes Trolle Petersen e adotterà il cognome Hjelmslev per distinguersi da Julius Petersen -- con il vecchio cognome resterà in matematica, tra l'altro, il teorema di Petersen--Morley. % lang-it
Sarà il padre del linguista Louis Hjelmslev. % lang-it
\index[persons]{Hjelmslev, Johannes}%
\index[persons]{Petersen, Julius}%

\begin{theorem}
[Hjelmsleva] % lang-pl
[di Hjelmslev] % lang-it
\index{twierdzenie!Hjelmsleva} % lang-pl
\index{teorema!di Hjelmslev} % lang-it
	Środki odcinków łączących odpowiadające sobie punkty dwóch przystających ciągów punktów (na płaszczyźnie euklidesowej lub Łobaczewskiego) są współliniowe albo pokrywają się. % lang-pl
	I punti medi dei segmenti che congiungono punti corrispondenti di due successioni congruenti di punti (nel piano euclideo o in quello di Lobachevskij) sono allineati oppure coincidono. % lang-it
\end{theorem}

Twierdzenie to znane będzie od dawna, ale dopiero w 1907 roku Hjelmslev odkryje, że należy ono do geometrii absolutnej i że ma liczne zastosowania w geometrii Łobaczewskiego; dowód: Eves \cite[s. 311--312]{eves1_1972}. % lang-pl
Il teorema sarà noto da tempo, ma solo nel 1907 Hjelmslev scoprirà che appartiene alla geometria assoluta e ha numerose applicazioni nella geometria di Lobachevskij; dimostrazione: Eves \cite[p. 311--312]{eves1_1972}. % lang-it

Ustalmy punkt $O$ i kąt ostry $\theta$. % lang-pl
Przekształcenie $T$ przypisuje punktowi $P$ taki punkt $P'$, że $\angle POP' = \theta$ i $\angle OP'P = \tfrac{\pi}{2}$, a następnie obraca $P'$ wokół $O$ o $-\theta$ \cite[s. 312]{eves1_1972}. % lang-pl
Fissiamo un punto $O$ e un angolo acuto $\theta$. % lang-it
La trasformazione $T$ associa a un punto $P$ il punto $P'$ tale che $\angle POP' = \theta$ e $\angle OP'P = \tfrac{\pi}{2}$, e poi ruota $P'$ intorno a $O$ di $-\theta$ \cite[p. 312]{eves1_1972}. % lang-it

\begin{proposition}
	Na płaszczyźnie euklidesowej $T$ jest jednokładnością o środku $O$ i skali $\cos\theta$, a więc przekształca płaszczyznę na siebie. % lang-pl
	Na płaszczyźnie Łobaczewskiego $T$ przekształca całą płaszczyznę na wnętrze koła o środku $O$ i promieniu $h$, gdzie $\Pi(h) = \theta$. % lang-pl
	Nel piano euclideo $T$ è l'omotetia di centro $O$ e rapporto $\cos\theta$, dunque manda il piano in sé. % lang-it
	Nel piano di Lobachevskij $T$ manda l'intero piano nell'interno del cerchio di centro $O$ e raggio $h$, dove $\Pi(h) = \theta$. % lang-it
\end{proposition}

Dowody: Eves \cite[s. 312--313]{eves1_1972}. % lang-pl
Eves zauważa też, że obraz jest podobny do modelu z sekcji \ref{sekcja_punkty_idealne}, i nie rozwija tego związku z braku miejsca \cite[s. 313]{eves1_1972}. % lang-pl
Dimostrazioni: Eves \cite[p. 312--313]{eves1_1972}. % lang-it
Eves nota anche che l'immagine somiglia al modello della sezione \ref{sekcja_punkty_idealne}, e non sviluppa questo legame per mancanza di spazio \cite[p. 313]{eves1_1972}. % lang-it

Hjelmslev będzie też jednym z autorów dowodów twierdzenia o jednorodności (zobacz sekcję \ref{sekcja_suma_katow}): według przypisu Hilberta, cytowanego przez Greenberga, dowód pierwszy poda Dehn, a późniejsze F.~Schur i Hjelmslev \cite[s. 208]{greenberg_2010}. % lang-pl
Hjelmslev sarà anche uno degli autori di dimostrazioni del teorema di uniformità (si veda la sezione \ref{sekcja_suma_katow}): secondo una nota di Hilbert, citata da Greenberg, la prima dimostrazione sarà di Dehn e quelle successive di F.~Schur e Hjelmslev \cite[p. 208]{greenberg_2010}. % lang-it

Zadania: Eves \cite[s. 314]{eves1_1972}, w tym wyprowadzenie z odwzorowania $T$ przenoszalności, symetrii i przechodniości równoległości. % lang-pl
Esercizi: Eves \cite[p. 314]{eves1_1972}, tra cui la deduzione dalla trasformazione $T$ della trasmissibilità, simmetria e transitività del parallelismo. % lang-it
